\documentclass[10pt,a4paper]{amsart}
\usepackage{amsmath,amssymb,amsthm,graphicx}
\usepackage{verbatim}
\newenvironment{reference}{\comment}{\endcomment}
%\newenvironment{reference}{}{}
\newenvironment{slogan}{\comment}{\endcomment}
\newenvironment{history}{\comment}{\endcomment}
% For commutative diagrams we use Xy-pic
\usepackage[all]{xy}

% We use 2cell for 2-commutative diagrams.
\xyoption{2cell}
% Theorem environments.
%
\theoremstyle{plain}
\newtheorem{theorem}[subsection]{Theorem}
\newtheorem{proposition}[subsection]{Proposition}
\newtheorem{lemma}[subsection]{Lemma}

\theoremstyle{definition}
\newtheorem{definition}[subsection]{Definition}
\newtheorem{example}[subsection]{Example}
\newtheorem{exercise}[subsection]{Exercise}
\newtheorem{situation}[subsection]{Situation}

\theoremstyle{remark}
\newtheorem{remark}[subsection]{Remark}
\newtheorem{remarks}[subsection]{Remarks}

\numberwithin{equation}{subsection}

% Macros
%
\def\lim{\mathop{\mathrm{lim}}\nolimits}
\def\colim{\mathop{\mathrm{colim}}\nolimits}
\def\Spec{\mathop{\mathrm{Spec}}}
\def\Hom{\mathop{\mathrm{Hom}}\nolimits}
\def\Ext{\mathop{\mathrm{Ext}}\nolimits}
\def\SheafHom{\mathop{\mathcal{H}\!\mathit{om}}\nolimits}
\def\SheafExt{\mathop{\mathcal{E}\!\mathit{xt}}\nolimits}
\def\Sch{\mathit{Sch}}
\def\Mor{\mathop{\mathrm{Mor}}\nolimits}
\def\Ob{\mathop{\mathrm{Ob}}\nolimits}
\def\Sh{\mathop{\mathit{Sh}}\nolimits}
\def\NL{\mathop{N\!L}\nolimits}
\def\CH{\mathop{\mathrm{CH}}\nolimits}
\def\proetale{{pro\text{-}\acute{e}tale}}
\def\etale{{\acute{e}tale}}
\def\QCoh{\mathit{QCoh}}
\def\Ker{\mathop{\mathrm{Ker}}}
\def\Im{\mathop{\mathrm{Im}}}
\def\Coker{\mathop{\mathrm{Coker}}}
\def\Coim{\mathop{\mathrm{Coim}}}

% Boxtimes
%
\DeclareMathSymbol{\boxtimes}{\mathbin}{AMSa}{"02}

%
% Macros for moduli stacks/spaces
%
\def\QCohstack{\mathcal{QC}\!\mathit{oh}}
\def\Cohstack{\mathcal{C}\!\mathit{oh}}
\def\Spacesstack{\mathcal{S}\!\mathit{paces}}
\def\Quotfunctor{\mathrm{Quot}}
\def\Hilbfunctor{\mathrm{Hilb}}
\def\Curvesstack{\mathcal{C}\!\mathit{urves}}
\def\Polarizedstack{\mathcal{P}\!\mathit{olarized}}
\def\Complexesstack{\mathcal{C}\!\mathit{omplexes}}
% \Pic is the operator that assigns to X its picard group, usage \Pic(X)
% \Picardstack_{X/B} denotes the Picard stack of X over B
% \Picardfunctor_{X/B} denotes the Picard functor of X over B
\def\Pic{\mathop{\mathrm{Pic}}\nolimits}
\def\Picardstack{\mathcal{P}\!\mathit{ic}}
\def\Picardfunctor{\mathrm{Pic}}
\def\Deformationcategory{\mathcal{D}\!\mathit{ef}}
\usepackage{cleveref}
\crefname{originallemma}{original lemma}{original lemmas}
\Crefname{originallemma}{Original lemma}{Original lemmas}
\makeatletter
\newlabel{lemma-gerbe-descent}{{lemma-gerbe-descent}{source}{}{sourcepointer-1}{}}
\newlabel{lemma-gerbe-descent@cref}{{[originallemma][0][]lemma-gerbe-descent}{[1][0][]source}{}{}{}}
\makeatother
\crefname{originallemma}{original lemma}{original lemmas}
\Crefname{originallemma}{Original lemma}{Original lemmas}
\makeatletter
\newlabel{lemma-gerbe-fppf}{{lemma-gerbe-fppf}{source}{}{sourcepointer-2}{}}
\newlabel{lemma-gerbe-fppf@cref}{{[originallemma][0][]lemma-gerbe-fppf}{[1][0][]source}{}{}{}}
\makeatother
\crefname{originallemma}{original lemma}{original lemmas}
\Crefname{originallemma}{Original lemma}{Original lemmas}
\makeatletter
\newlabel{lemma-gerbe-isom-fppf}{{lemma-gerbe-isom-fppf}{source}{}{sourcepointer-3}{}}
\newlabel{lemma-gerbe-isom-fppf@cref}{{[originallemma][0][]lemma-gerbe-isom-fppf}{[1][0][]source}{}{}{}}
\makeatother
\crefname{originallemma}{original lemma}{original lemmas}
\Crefname{originallemma}{Original lemma}{Original lemmas}
\makeatletter
\newlabel{lemma-gerbe-over-iso-classes}{{lemma-gerbe-over-iso-classes}{source}{}{sourcepointer-4}{}}
\newlabel{lemma-gerbe-over-iso-classes@cref}{{[originallemma][0][]lemma-gerbe-over-iso-classes}{[1][0][]source}{}{}{}}
\makeatother
\crefname{originallemma}{original lemma}{original lemmas}
\Crefname{originallemma}{Original lemma}{Original lemmas}
\makeatletter
\newlabel{lemma-surjective-flat-locally-finite-presentation}{{lemma-surjective-flat-locally-finite-presentation}{source}{}{sourcepointer-5}{}}
\newlabel{lemma-surjective-flat-locally-finite-presentation@cref}{{[originallemma][0][]lemma-surjective-flat-locally-finite-presentation}{[1][0][]source}{}{}{}}
\makeatother
\crefname{originalproposition}{original proposition}{original propositions}
\Crefname{originalproposition}{Original proposition}{Original propositions}
\makeatletter
\newlabel{proposition-when-gerbe}{{proposition-when-gerbe}{source}{}{sourcepointer-6}{}}
\newlabel{proposition-when-gerbe@cref}{{[originalproposition][0][]proposition-when-gerbe}{[1][0][]source}{}{}{}}
\makeatother
\title{A morphism is a gerbe exactly when both it and its diagonal are surjective, flat and of finite presentation}
\author{The Stacks Project Authors}
\date{stacks-project:stacks-morphisms.tex:local-snapshot-20261012; GNU FDL 1.2 or later}
\begin{document}
\maketitle
\noindent\textit{Editorial note.} All delivered mathematical text is the author's own, copied from the pinned Stacks Project chapter source sources/stacks-algebraic-stacks-morphisms/pinned-original-source.tex; no statement or proof step was rewritten, reordered or removed.\par\medskip
\noindent\textit{Editorial note.} Preamble provenance: the Stacks Project keeps one single shared preamble for every chapter; this card takes it directly from the staged sources/stacks-algebraic-stacks-morphisms/preamble.tex as three verbatim spans (lines 8-12, 14-18 and 154-217), so no derived file was written and nothing was added or redefined. The wrapper class and the base packages are supplied by the assembler.\par\medskip
\noindent\textit{Editorial note.} Bibliography provenance: neither the counted statement nor its proof contains a citation (this lane rescans the whole span before writing the card), so the card carries no bibliography and no entry was reconstructed.\par\medskip
\noindent\textit{Editorial note.} Selection provenance: this card was opened under the file-based control contract (CONTROL.md), with the mechanical pre-card gate run on its statement and proof interval before the card was written.\par\medskip
\noindent\textit{Allowed standard prerequisites (editorial).}\begin{itemize}
\item The two clauses of the gerbe condition through local existence of objects and of isomorphisms. Used inside the counted proof (source lines 6532-6545). Supplies the local objects and isomorphisms from the flatness and finite presentation hypotheses.
\item The statement that a morphism with surjective flat finitely presented diagonal and structure morphism has a flat finitely presented inertia. Invoked inside the counted proof (source lines 6546-6552). Connects the diagonal to the inertia.
\item The criterion for flatness and local finite presentation along a flat finitely presented cover. Invoked inside the counted proof (source lines 6553-6560). Allows the gerbe condition to be checked after base change.
\end{itemize}\medskip
\ifdefined\hypertarget\hypertarget{sourcepointer-1}{}\fi
\noindent\textit{Source reference.} Original lemma lemma-gerbe-descent: sources/stacks-algebraic-stacks-morphisms/pinned-original-source.tex, source lines 6161--6161. Navigation pointer only: the counted text refers to the Stacks item lemma-gerbe-descent, which is an existing labeled result of this same chapter and is not part of this card. The pointer carries the target's own label as its number and its definition line as its span; its content is not reproduced here.\par\medskip
\ifdefined\hypertarget\hypertarget{sourcepointer-2}{}\fi
\noindent\textit{Source reference.} Original lemma lemma-gerbe-fppf: sources/stacks-algebraic-stacks-morphisms/pinned-original-source.tex, source lines 6287--6287. Navigation pointer only: the counted text refers to the Stacks item lemma-gerbe-fppf, which is an existing labeled result of this same chapter and is not part of this card. The pointer carries the target's own label as its number and its definition line as its span; its content is not reproduced here.\par\medskip
\ifdefined\hypertarget\hypertarget{sourcepointer-3}{}\fi
\noindent\textit{Source reference.} Original lemma lemma-gerbe-isom-fppf: sources/stacks-algebraic-stacks-morphisms/pinned-original-source.tex, source lines 6429--6429. Navigation pointer only: the counted text refers to the Stacks item lemma-gerbe-isom-fppf, which is an existing labeled result of this same chapter and is not part of this card. The pointer carries the target's own label as its number and its definition line as its span; its content is not reproduced here.\par\medskip
\ifdefined\hypertarget\hypertarget{sourcepointer-4}{}\fi
\noindent\textit{Source reference.} Original lemma lemma-gerbe-over-iso-classes: sources/stacks-algebraic-stacks-morphisms/pinned-original-source.tex, source lines 6091--6091. Navigation pointer only: the counted text refers to the Stacks item lemma-gerbe-over-iso-classes, which is an existing labeled result of this same chapter and is not part of this card. The pointer carries the target's own label as its number and its definition line as its span; its content is not reproduced here.\par\medskip
\ifdefined\hypertarget\hypertarget{sourcepointer-5}{}\fi
\noindent\textit{Source reference.} Original lemma lemma-surjective-flat-locally-finite-presentation: sources/stacks-algebraic-stacks-morphisms/pinned-original-source.tex, source lines 5877--5877. Navigation pointer only: the counted text refers to the Stacks item lemma-surjective-flat-locally-finite-presentation, which is an existing labeled result of this same chapter and is not part of this card. The pointer carries the target's own label as its number and its definition line as its span; its content is not reproduced here.\par\medskip
\ifdefined\hypertarget\hypertarget{sourcepointer-6}{}\fi
\noindent\textit{Source reference.} Original proposition proposition-when-gerbe: sources/stacks-algebraic-stacks-morphisms/pinned-original-source.tex, source lines 6319--6319. Navigation pointer only: the counted text refers to the Stacks item proposition-when-gerbe, which is an existing labeled result of this same chapter and is not part of this card. The pointer carries the target's own label as its number and its definition line as its span; its content is not reproduced here.\par\medskip
\begin{proposition}
\label{proposition-when-gerbe-over}
Let $f : \mathcal{X} \to \mathcal{Y}$ be a morphism of algebraic stacks.
The following are equivalent
\begin{enumerate}
\item $\mathcal{X}$ is a gerbe over $\mathcal{Y}$, and
\item $f : \mathcal{X} \to \mathcal{Y}$ and
$\Delta : \mathcal{X} \to \mathcal{X} \times_\mathcal{Y} \mathcal{X}$
are surjective, flat, and locally of finite presentation.
\end{enumerate}
\end{proposition}

\begin{proof}
The implication (1) $\Rightarrow$ (2) follows from
Lemmas \ref{lemma-gerbe-fppf} and \ref{lemma-gerbe-isom-fppf}.

\medskip\noindent
Assume (2). It suffices to prove (1) for the base change of $f$
by a surjective, flat, and locally finitely presented morphism
$\mathcal{Y}' \to \mathcal{Y}$, see
Lemma \ref{lemma-gerbe-descent} (note that the base change
of the diagonal of $f$ is the diagonal of the base change). 
Thus we may assume $\mathcal{Y}$ is a scheme $Y$.
In this case $\mathcal{I}_\mathcal{X} \to \mathcal{X}$
is a base change of $\Delta$ and
we conclude that $\mathcal{X}$ is a gerbe
by Proposition \ref{proposition-when-gerbe}.
We still have to show that $\mathcal{X}$ is a gerbe over $Y$.
Let $\mathcal{X} \to X$ be the morphism of
Lemma \ref{lemma-gerbe-over-iso-classes}
turning $\mathcal{X}$ into a gerbe over the algebraic space $X$
classifying isomorphism classes of objects of $\mathcal{X}$.
It is clear that $f : \mathcal{X} \to Y$ factors as
$\mathcal{X} \to X \to Y$. Since $f$ is surjective, flat, and
locally of finite presentation, we conclude that $X \to Y$ is
surjective as a map of fppf sheaves (for example use
Lemma \ref{lemma-surjective-flat-locally-finite-presentation}).
On the other hand, $X \to Y$ is injective too: for any scheme
$T$ and any two $T$-valued points $x_1, x_2$ of $X$ which map to
the same point of $Y$, we can first fppf locally on $T$
lift $x_1, x_2$ to objects $\xi_1, \xi_2$ of $\mathcal{X}$ over $T$
and second deduce that $\xi_1$ and $\xi_2$ are fppf locally isomorphic
by our assumption that
$\Delta : \mathcal{X} \to \mathcal{X} \times_Y \mathcal{X}$
is surjective, flat, and locally of finite presentation.
Whence $x_1 = x_2$ by construction of $X$.
Thus $X = Y$ and the proof is complete.
\end{proof}

\end{document}
